El presente capítulo expone la evaluación piloto del prototipo desarrollado, en correspondencia con el \textbf{Objetivo Específico 5}, orientado a:

\begin{quote}
Evaluar la efectividad de los módulos en términos de comprensión del contenido, participación, motivación, retención de información y aplicación práctica, mediante pruebas piloto y retroalimentación de los estudiantes, identificando áreas de mejora.
\end{quote}

Para ello se aplicaron tres encuestas dirigidas a estudiantes de Ingeniería de Sistemas, cuyos resultados permiten analizar los criterios de evaluación definidos en este objetivo.
\\
\\
Con el propósito de mantener el capítulo enfocado en el análisis, se presenta aquí una síntesis estructurada de los hallazgos, mientras que la evidencia detallada (preguntas, respuestas completas y gráficos por pregunta) se incluye en el \hyperref[anexo:encuestas]{Anexo I}.

\section{Justificación de las pruebas}

La aplicación de pruebas a los estudiantes es fundamental para determinar en qué medida el prototipo contribuye al desarrollo de las habilidades blandas trabajadas en este proyecto. En particular, las encuestas permiten:

\begin{itemize}
    \item Obtener información directa sobre la experiencia de los estudiantes con los módulos.
    \item Identificar fortalezas, dificultades y oportunidades de mejora desde la perspectiva de los estudiantes.
    \item Validar el diseño del prototipo y comprender su impacto formativo.
    \item Orientar ajustes futuros para optimizar la experiencia de aprendizaje.
\end{itemize}

\section{Metodología de aplicación}

El proceso de evaluación del prototipo se desarrolló con dos poblaciones objetivo diferenciadas, con el fin de obtener perspectivas complementarias sobre la usabilidad de la plataforma y la efectividad percibida de los módulos formativos.

\subsection*{Población 1: Evaluación de usabilidad (Heurísticas de Nielsen)}
\begin{itemize}
    \item \textbf{Población objetivo:} Egresados y profesionales de Ingeniería de Sistemas.
    
    \item \textbf{Justificación:} Esta población cuenta con experiencia técnica y criterio profesional para evaluar la interfaz con mayor profundidad, lo que facilita identificar problemas de usabilidad, consistencia y diseño que pueden pasar desapercibidos para los usuarios finales.
    
    \item \textbf{Enfoque:} Aplicación de una encuesta construida a partir de las heurísticas de Nielsen, adaptadas al contexto educativo.
\end{itemize}

\subsection*{Población 2: Evaluación de efectividad percibida (usuarios finales)}
\begin{itemize}
    \item \textbf{Población objetivo:} Estudiantes de Ingeniería de Sistemas, usuarios finales del prototipo.
    
    \item \textbf{Justificación:} La participación de los estudiantes es fundamental, puesto que el prototipo fue diseñado para apoyar su proceso de aprendizaje y solo ellos pueden ofrecer una valoración basada en su propia experiencia de uso.
    
    \item \textbf{Enfoque:} Aplicación de dos encuestas centradas en los módulos específicos Aprender a Aprender y Escucha Activa.
\end{itemize}

\section{Descripción de las encuestas aplicadas}

A continuación se describen los tres instrumentos utilizados en la evaluación del prototipo. Todas las encuestas combinaron preguntas cerradas tipo Likert con preguntas abiertas, lo que permitió obtener tanto valoraciones cuantitativas como observaciones cualitativas sobre el funcionamiento de los módulos y la experiencia de uso de la plataforma.


\begin{itemize}
    \item \textbf{Encuesta 1: Evaluación del módulo Aprender a Aprender.}  
    Instrumento orientado a identificar cómo los estudiantes comprendieron los componentes centrales del módulo (metas, objetivos, tareas y roadmap), así como la utilidad percibida de las herramientas implementadas (tablero Kanban, editor visual de roadmaps y autoevaluaciones).
    La encuesta recoge percepciones sobre claridad, participación, motivación, retención de conceptos clave y aplicación práctica, complementadas con preguntas abiertas que permiten identificar experiencias, dificultades y sugerencias.


    \item \textbf{Encuesta 2: Evaluación del módulo Escucha Activa.}  
    Instrumento diseñado para valorar la experiencia de los estudiantes con las actividades auditivas del módulo y su contribución al desarrollo de habilidades de escucha activa.
    Evalúa la comprensión de los principios trabajados, la utilidad de la retroalimentación, la capacidad para recordar información, y la aplicabilidad de lo aprendido en contextos reales.
    Incluye además preguntas abiertas que permiten identificar audios más útiles, dificultades frecuentes y recomendaciones de mejora.


    \item \textbf{Encuesta 3: Evaluación de usabilidad de la plataforma.}  
    
    Instrumento dirigido a valorar la usabilidad del prototipo tomando como referencia las heurísticas de Nielsen. La encuesta incluye ítems destinados a medir la claridad del estado del sistema, la consistencia de la interfaz, la facilidad de navegación, la prevención de errores y la comprensión de mensajes e indicadores visuales. 
    
    Complementariamente, incorpora preguntas abiertas orientadas a identificar elementos confusos, dificultades de uso y sugerencias de mejora relacionadas con la experiencia general de usuario.

\end{itemize}

\section{Proceso de aplicación}

La recolección de información se desarrolló mediante un procedimiento mixto que combinó momentos de aplicación guiada con participación asincrónica. Cada encuesta fue distribuida según la población objetivo correspondiente, junto con una breve explicación del propósito del estudio y las instrucciones para su diligenciamiento.
\\
\\
En el caso de la evaluación de usabilidad, el enlace fue compartido con egresados y profesionales de Ingeniería de Sistemas, quienes exploraron libremente el prototipo antes de completar la encuesta basada en las heurísticas de Nielsen.
\\
\\
Para la evaluación de los módulos \textit{Aprender a Aprender} y \textit{Escucha Activa}, se contó con la colaboración de un docente del programa, quien facilitó un espacio presencial de aproximadamente 30 minutos en el que los estudiantes interactuaron con el prototipo y posteriormente diligenciaron las encuestas correspondientes. Este espacio permitió asegurar que todos los participantes tuvieran un primer acercamiento guiado al funcionamiento del sistema.
\\
\\
Adicionalmente, se ofreció la posibilidad de responder las encuestas de manera asincrónica para aquellos estudiantes que no lograron completar el formulario durante el espacio asignado. Con ello se garantizó la participación de un mayor número de estudiantes y se respetaron sus tiempos individuales.


\section{Limitaciones del proceso de evaluación}

Si bien las encuestas aplicadas permitieron obtener información valiosa sobre la experiencia de uso y la efectividad percibida del prototipo, el proceso de evaluación presenta algunas limitaciones que deben considerarse al interpretar los resultados.

\begin{itemize}
    \item \textbf{Alcance exploratorio:} La evaluación corresponde a una fase temprana de validación del prototipo, lo que implica que los resultados se enfocan en percepciones y no permiten establecer conclusiones causales sobre el desarrollo de habilidades.

    \item \textbf{Tamaño y composición de la muestra:} La participación estuvo limitada a estudiantes y egresados disponibles durante el periodo de aplicación, lo que reduce la posibilidad de generalizar los resultados a toda la población estudiantil.

    \item \textbf{Tiempo de interacción con el prototipo:} Aunque se contó con un espacio guiado de 30 minutos, este tiempo puede no haber sido suficiente para explorar en profundidad todas las funcionalidades, especialmente en el módulo \textit{Aprender a Aprender}, que involucra procesos más prolongados.

    \item \textbf{Autorreporte:} Al basarse en encuestas, los resultados dependen de la autorreflexión y percepción subjetiva de los participantes, lo cual puede introducir sesgos asociados a expectativas, motivación u otros factores personales.

    \item \textbf{Necesidad de experiencia sostenida:} Algunos componentes del prototipo, como la construcción de hábitos o la planificación a largo plazo, requieren un uso continuo que no pudo evaluarse en el marco de una sesión breve.

    \item \textbf{Ausencia de grupo control:} No se incluyó un grupo que no utilizara el prototipo, por lo que no fue posible realizar comparaciones.
\end{itemize}

Estas limitaciones no invalidan los hallazgos obtenidos; más bien permiten contextualizar su alcance y abren la puerta a futuras investigaciones y evaluaciones más amplias.






